% Einheit 4 von 4 – Zuordnungstypen erkennen und anwenden (bank/zuordnungen/e4.jsonl, zusammenbau v0.1)
%% TODO Zweigzeile Teil 1: was hier gelernt wird (Satz in Schülersprache aus dem Einheitstitel) – Einheit 4
\einheitenkopf[e4]{Einheit 4 von 4 · Zuordnungstypen erkennen und anwenden}
\zweigzeile{ab Kl. 7 · P10 oft}
\setcounter{aufgabe}{22}

%% TODO Titel: Ich-kann-Satz fehlt in der Bank (Platzhalter: Kettenname) – Nr. 23
%% TODO Anweisung: ein Satz über der ersten Teilaufgabe fehlt in der Bank – Nr. 23
\begin{aufgabe}{Nullwert prüfen}
\begin{teile}
\teil Taxi: 4 € Grundgebühr und 2 € je km. Kostet eine Fahrt von 0 km auch 0 €? \\ \kreuz{ja} \\ \kreuz{nein}
\end{teile}
\end{aufgabe}

%% TODO \verfahren{…}: Verfahrensüberschrift in Sachsprache fehlt in der Bank (2.3 b)
%% TODO Titel: Ich-kann-Satz fehlt in der Bank (Platzhalter: Kettenname) – Nr. 24
%% TODO Anweisung: ein Satz über der ersten Teilaufgabe fehlt in der Bank – Nr. 24
\begin{aufgabe}{Erkennen}
%% TODO \rechenplatz{n}: Schrittzahl des Grundfalls fehlt in der Bank (nur bei fester Schreibform, 2.3 b)
\begin{teile}
\teil Je mehr Hefte du kaufst, desto … bezahlst du. \\ \kreuz{mehr} \\ \kreuz{weniger} \\ \kreuz{nicht vorhersagbar}
\teil Ist die Zuordnung in der Tabelle proportional? Prüfe die Quotienten y : x. \\ \kreuz{proportional} \\ \kreuz{nicht proportional}

\wertetabelle[6,12,15]{x}{y}{2,4,5}
\teil Ist die Zuordnung in der Tabelle antiproportional? Prüfe die Produkte x · y. \\ \kreuz{antiproportional} \\ \kreuz{nicht antiproportional}

\wertetabelle[9,6,3]{x}{y}{2,3,6}
\teil Fahrradverleih: 3 € Grundgebühr und 2 € je Stunde. Ist Mietdauer → Preis proportional? \\ \kreuz{proportional} \\ \kreuz{nicht proportional}
\teil Welcher Zuordnungstyp gehört zum Graphen? \\ \kreuz{proportional} \\ \kreuz{antiproportional} \\ \kreuz{keins von beiden}

\begin{ksys}[xmin=0,xmax=5,ymin=0,ymax=8]\gerade{1.5}{0}{}\end{ksys}
\teil Welcher Zuordnungstyp liegt in der Tabelle vor? Prüfe erst die Quotienten, dann die Produkte. \\ \kreuz{proportional} \\ \kreuz{antiproportional} \\ \kreuz{keins von beiden}

\wertetabelle[20,10,5]{x}{y}{2,4,8}
\teil Bei einem Paketdienst kosten 2 kg 8 € und 4 kg 14 €. Ist Gewicht → Preis proportional? Entscheide und begründe.
\end{teile}
\end{aufgabe}

%% TODO \verfahren{…}: Verfahrensüberschrift in Sachsprache fehlt in der Bank (2.3 b)
%% TODO Titel: Ich-kann-Satz fehlt in der Bank (Platzhalter: Kettenname) – Nr. 25
%% TODO Anweisung: ein Satz über der ersten Teilaufgabe fehlt in der Bank – Nr. 25
\begin{aufgabe}{Rate}
%% TODO \rechenplatz{n}: Schrittzahl des Grundfalls fehlt in der Bank (nur bei fester Schreibform, 2.3 b)
\begin{teile}
\teil Auf dem Tankbeleg steht: 40 l, 1,79 € je Liter, gesamt 71,60 €. Welche Angabe ist die Rate?
\teil Äpfel kosten 2,50 € je kg. Was kosten 4 kg? \leerfeld[€]
\teil Kartoffeln kosten 1,50 € je kg. Wie viele Kilogramm bekommt man für 6 €? \leerfeld[kg]
\teil Ein Radfahrer fährt 45 km in 3 Stunden. Wie schnell fährt er im Schnitt? \leerfeld km/h
\teil Ein Bus fährt 150 km mit durchschnittlich 50 km/h. Wie lange dauert die Fahrt? \leerfeld[h]
\teil Ein Lieferwagen verbraucht 8 l Diesel auf 100 km und fährt im Jahr 30\,000 km. Ein Liter Diesel kostet 162,5 ct. Berechne die Dieselkosten für ein Jahr in Euro. \hfill (P10 2014 OS) \\ \leerfeld[€]
\teil Eine Wandergruppe geht mit 4 km/h. Sie ist um 9:50 Uhr an einer Hütte, macht dort 30 Minuten Pause und geht dann 3\,000 m bis zum See. Wann kommt sie am See an? \hfill (P10 2014 OS) \\ um \leerfeld Uhr
\end{teile}
\end{aufgabe}

%% TODO Titel: Ich-kann-Satz fehlt in der Bank (Platzhalter: Kettenname) – Nr. 26
%% TODO Anweisung: ein Satz über der ersten Teilaufgabe fehlt in der Bank – Nr. 26
\begin{aufgabe}{Gegenbeispiele (Alter und Größe, Tarif mit Grundgebühr)}
\begin{teile}
\teil Nenne ein Beispiel für eine Zuordnung „je mehr …, desto mehr …“, die nicht proportional ist, und begründe kurz.
\end{teile}
\end{aufgabe}

%% TODO Titel: Ich-kann-Satz fehlt in der Bank (Platzhalter: Kettenname) – Nr. 27
%% TODO Anweisung: ein Satz über der ersten Teilaufgabe fehlt in der Bank – Nr. 27
\begin{aufgabe}{Erkennen – Fehler finden, Begründen, Darstellungswechsel, Anwendung}
\begin{teile}
\teil Ein Auto fährt 90 km in 1 Stunde. Lukas rechnet: „Für 180 km braucht es eine halbe Stunde, denn die Strecke ist doppelt so lang.“ Finde den Fehler und rechne richtig.
\teil Warum ist die Zuordnung Anzahl Personen → Preis pro Person für eine Pizza nicht proportional? Begründe.
\teil Ein Brunnen liefert 5 l Wasser je Minute. Zeichne den Graphen Zeit in min → Wasser in l für 0 bis 4 Minuten und nenne den Zuordnungstyp.

\begin{ksys}[xmin=0,xmax=5,ymin=0,ymax=20,xlabel=Zeit in min,ylabel=Wasser in l]\end{ksys}
\teil Familie Yilmaz fährt 280 km in den Urlaub. Sie fährt im Schnitt 80 km/h und will um 14:00 Uhr ankommen. Wann muss sie spätestens losfahren?
\end{teile}
\end{aufgabe}
